% from ega4-2-fr.tex
\oldpage[IV-2]{35}
(homomorphisme déterminé de façon unique (lorsqu'il existe) par cette
propriété), alors le diagramme
\[
\xymatrix{
\mathcal F_1 \ar[r]^u \ar[d] & \mathcal F_2 \ar[d] \\
i_*(\mathcal G'_1) \ar[r]_{i_*(v)} & i_*(\mathcal G'_2)
}
\]
est commutatif, et il y a par suite un seul homomorphisme
$w:\overline{\mathcal G'_1}\to\overline{\mathcal G'_2}$ rendant commutatif
le diagramme
\[
\xymatrix{
\mathcal F_1 \ar[r]^u \ar[d] & \mathcal F_2 \ar[d] \\
\overline{\mathcal G'_1} \ar[r]_w & \overline{\mathcal G'_2}
}
\]
\end{env}

\begin{proposition}[2.8.4]
\label{IV.2.8.4-fr}
Les hypothèses sur $Y$ étant celles de (2.8.1), soient $X_1$, $X_2$ deux
$Y$-préschémas, $\mathcal F_i$ un $\mathcal O_{X_i}$-Module
quasi-cohérent, $\mathcal G'_i$ un $\mathcal O_{(X_i)_\eta}$-Module quotient
de $(\mathcal F_i)_\eta$ $(i=1,2)$. Alors on a
\end{proposition}

\begin{env}[2.8.4.1]
\label{IV.2.8.4.1-fr}
\[
\overline{\mathcal G'_1}\otimes_{\mathcal O_Y}\overline{\mathcal G'_2}
=\overline{\mathcal G'_1\otimes_{k(\eta)}\mathcal G'_2}.
\]
\end{env}

\begin{proof}
En effet, posons $X=X_1\times_Y X_2$; le premier membre de (2.8.3.1) est
un $\mathcal O_X$-Module quasi-cohérent qui est $Y$-plat
$(0_{\mathrm I},6.2.1)$, dont l'image réciproque dans $X_\eta$ est
$\mathcal G'_1\otimes_{k(\eta)}\mathcal G'_2$ (I, 9.1.5) et qui est un
quotient de $\mathcal F_1\otimes_{\mathcal O_Y}\mathcal F_2$; la conclusion
résulte donc de la propriété d'unicité de (2.8.1).
\end{proof}

\begin{proposition}[2.8.5]
\label{IV.2.8.5-fr}
Les hypothèses sur $X$ et $Y$ étant celles de (2.8.1), soit $Z'$ un
sous-préschéma fermé de $X_\eta$. Il existe alors un unique sous-préschéma
fermé $\overline{Z'}$ de $X$ qui soit $Y$-plat et tel que
$i^{-1}(\overline{Z'})=Z'$.
\end{proposition}

\begin{proof}
Si $\mathcal J'$ est l'Idéal quasi-cohérent de $\mathcal O_{X_\eta}$
définissant $Z'$, il suffit d'appliquer (2.8.1) au cas où
$\mathcal F=\mathcal O_X$ et
$\mathcal G'=\mathcal O_{X_\eta}/\mathcal J'$; si
$\overline{\mathcal G'}=\mathcal O_X/\mathcal J$, on a en effet
$\mathcal J'=i^*(\mathcal J)\mathcal O_{X_\eta}$, donc
$i^{-1}(\overline{Z'})=Z'$ (I, 4.4.5).
\end{proof}

On notera que le préschéma $\overline{Z'}$ est l'\emph{image fermée} de $Z'$
par le morphisme composé
\[
Z'\longrightarrow X_\eta\xrightarrow{i}X,
\]
où la première flèche est l'injection canonique (I, 9.5.3); son espace
sous-jacent est l'adhérence \emph{dans $X$} de $Z'$ (I, 9.5.4), ce qui
justifie la notation adoptée. On dit encore que $\overline{Z'}$ est le
sous-préschéma de $X$ \emph{adhérence} de $Z'$.

\begin{corollary}[2.8.6]
\label{IV.2.8.6-fr}
Soient $X_1$, $X_2$ deux $Y$-préschémas, $Z'_i$ un sous-préschéma fermé de
$(X_i)_\eta$ $(i=1,2)$. Alors on a
\end{corollary}

\begin{env}[2.8.6.1]
\label{IV.2.8.6.1-fr}
\[
\overline{Z'_1}\times_Y\overline{Z'_2}
=\overline{Z'_1\times_{k(\eta)}Z'_2}.
\]
\end{env}

\begin{proof}
Cela résulte de (2.8.4) et (2.8.5).
\end{proof}

